\section*{Capítulo \ref{ch:b1:lab-techniques-1} --- Técnicas de laboratório I: segurança, medição e separação}

\begin{solution}{exo:b1:lab-techniques-1:1}
\emph{Perigo}: a H225 (líquido altamente inflamável, categoria 2) o exige,
e o rótulo traz a palavra mais grave. A H225 é um \omterm{def:b1:lab-techniques-1:hazard}{perigo} físico; a H319
(irritação ocular grave) e a H336 (sonolência ou vertigem) são \omterm{def:b1:lab-techniques-1:hazard}{perigos}
para a saúde. Precauções: nenhuma chama ou faísca por perto, um frasco
fechado; óculos; trabalhar em local ventilado ou em capela.
% ledger: ghs:acetone, hstat:H225, hstat:H319, hstat:H336
\end{solution}

\begin{solution}{exo:b1:lab-techniques-1:2}
(a) Um \omterm{def:b1:lab-techniques-1:hazard}{perigo}, uma propriedade do ácido. (b) Um \omterm{def:b1:lab-techniques-1:hazard}{risco}, pequeno porque a
exposição é pequena. (c) Um \omterm{def:b1:lab-techniques-1:hazard}{perigo}. (d) Um \omterm{def:b1:lab-techniques-1:hazard}{risco}: o mesmo \omterm{def:b1:lab-techniques-1:hazard}{perigo}, maior
exposição (vapor, respingos) a quente e em recipiente aberto.
\end{solution}

\begin{solution}{exo:b1:lab-techniques-1:3}
$R_f = 2.1/6.0 = 0.35$ e $4.2/6.0 = 0.70$. A segunda mancha tem o mesmo
$R_f$ que a referência: a amostra pode contê-la (é coerente), mas $R_f$
iguais com um único eluente não provam a identidade; um segundo eluente,
ou um depósito conjunto da amostra e da referência que continue sendo uma
única mancha, reforça a conclusão. A mancha a 0.35 é com certeza outra espécie.
\end{solution}

\begin{solution}{exo:b1:lab-techniques-1:4}
Semiamplitude \qty{0.05}{mL}, retangular: $u = 0.05/\sqrt3 = \qty{0.029}{mL}$.
Para uma diferença de duas leituras independentes,
$u = \sqrt{2} \times 0.029 = \qty{0.041}{mL}$.
\end{solution}

\begin{solution}{exo:b1:lab-techniques-1:5}
Média \qty{0.25100}{g}; desvios $+2$, $-2$, $+5$, 0, $-5$ (em
\qty{e-4}{g}), soma dos quadrados \qty{58e-8}{g^2}, $s = \sqrt{58 \times
10^{-8}/4} = \qty{3.8e-4}{g}$; $u = s/\sqrt5 = \qty{1.7e-4}{g}$;
$U = \qty{3.4e-4}{g}$: $m = (0.25100 \pm 0.00034)$~g, $k = 2$.
\end{solution}

\begin{solution}{exo:b1:lab-techniques-1:6}
Uma porção: $f = 1/(1 + 3 \times 30/50) = 0.357$, \qty{64.3}{\percent}
extraídos. Duas de \qty{15}{mL}: $f = 1.9^{-2} = 0.277$,
\qty{72.3}{\percent}. Três de \qty{10}{mL}: $f = 1.6^{-3} = 0.244$,
\qty{75.6}{\percent}.
\end{solution}

\begin{solution}{exo:b1:lab-techniques-1:7}
$z = 0.0012/\sqrt{0.0006^2 + 0.0002^2} = 0.0012/0.00063 = 1.9 \le 2$:
compatíveis, por pouco. Com \qty{0.0004}{mol/L}: $z = 0.0012/0.00045 = 2.7$,
não compatíveis: a \omterm{def:b1:titration-methods:titration}{titulação} mais precisa revela uma diferença real (um
erro na preparação, ou na \omterm{def:b1:titration-methods:titration}{titulação}).
\end{solution}

\begin{solution}{exo:b1:lab-techniques-1:8}
Q: pouco solúvel a frio, muito solúvel a quente. P dissolve o sólido
quase tão mal a quente quanto a frio (nada cristaliza no resfriamento); R
mantém dissolvido demais a frio (grandes perdas). Com Q, \qty{5.0}{g}
exigem pelo menos $5.0/12 \times 100 = \qty{42}{mL}$ de solvente em
ebulição, que mantêm $0.3 \times 0.417 = \qty{0.13}{g}$ dissolvidos a
frio: no máximo \qty{4.9}{g}, \qty{97.5}{\percent}, podem ser recuperados
(antes de qualquer perda no filtro).
\end{solution}

\begin{solution}{exo:b1:lab-techniques-1:9}
Etanol ($n_D = 1.3611$ a \qty{20}{\celsius}); a água (1.333) e o
ciclo-hexano (1.4266) estão excluídos. O índice varia com a temperatura,
de modo que um valor nada significa sem ela. Um valor de 1.350, entre a
água e o etanol, sugere uma mistura, como etanol contendo água.
% ledger: nD:water, nD:ethanol, nD:cyclohexane
\end{solution}

\begin{solution}{exo:b1:lab-techniques-1:10}
Com $t = KV/V_a$ e $n$ porções, $\ln f_n = -n\ln(1 + t/n)$; quando
$n \to \infty$, $n \ln(1 + t/n) \to t$ (pois $\ln(1 + \varepsilon) \sim
\varepsilon$), e $f_n$ decresce até $\mathrm{e}^{-t}$ (\cref{prop:b1:lab-techniques-1:multiple-extractions}).
Aqui $t = 3 \times 30/50 = 1.8$, $\mathrm{e}^{-1.8} = 0.165$: no máximo
\qty{83.5}{\percent}. Três porções já dão \qty{75.6}{\percent},
\qty{91}{\percent} do limite; atingir \qty{99}{\percent} dele exige 32
porções de menos de \qty{1}{mL}: além de duas ou três porções, o ganho
não compensa o trabalho.
\end{solution}

\begin{solution}{exo:b1:lab-techniques-1:11}
Com $t = x - x_0$: $\int_{-a}^{a} (a - |t|)/a^2 \,\mathrm{d}t =
2 \int_0^a (a - t)/a^2 \,\mathrm{d}t = 2 (a^2/2)/a^2 = 1$. A média é
$x_0$ por simetria; a variância é $2 \int_0^a t^2 (a - t)/a^2
\,\mathrm{d}t = \frac{2}{a^2}\left(\frac{a^4}{3} - \frac{a^4}{4}\right) =
\frac{a^2}{6}$, logo $u = a/\sqrt6 = 0.41\,a$, contra $a/\sqrt3 = 0.58\,a$
no caso retangular: saber que o meio é mais provável reduz a
incerteza.
\end{solution}

\begin{solution}{exo:b1:lab-techniques-1:12}
$c = m/(MV) = 0.25100/(204.2 \times 0.10000) = \qty{0.012292}{mol/L}$.
\omterm{def:b1:lab-techniques-1:uncertainty}{Incertezas relativas}: massa $1.7 \times 10^{-4}/0.2510 = 6.8 \times
10^{-4}$, volume $0.06/100.00 = 6.0 \times 10^{-4}$; combinada
$\sqrt{(6.8^2 + 6.0^2)} \times 10^{-4} = 9.1 \times 10^{-4}$, isto é,
\qty{0.091}{\percent}; $u(c) = \qty{1.1e-5}{mol/L}$, $U = \qty{2.2e-5}{mol/L}$:
$c = (0.012292 \pm 0.000022)$~mol/L, $k = 2$. A pesagem domina, por
pouco; as duas fontes são da mesma ordem.
\end{solution}

\begin{solution}{pb:b1:lab-techniques-1:1}
\textbf{1.} GHS02 inflamável, GHS05 corrosivo, GHS06 toxicidade aguda,
GHS07 nocivo ou irritante.
\textbf{2.} Ácido salicílico: \emph{Perigo}, por causa da H318 (lesões
oculares graves, categoria 1). Aspirina: \emph{Atenção} (só H302).
\textbf{3.} O \omterm{def:b1:lab-techniques-1:hazard}{perigo} é fixo, a exposição não: um volume pequeno, medido e
usado em capela, com luvas, óculos e jaleco, o frasco fechado de
imediato, nenhuma chama por perto.
\textbf{4.} H314 (provoca queimadura severa à pele e dano aos olhos):
óculos e luvas. H332 (nocivo se inalado) e H226 (líquido e vapores
inflamáveis): a capela e a ausência de chamas.
\textbf{5.} Enxaguar cuidadosamente com água durante vários minutos,
removendo as lentes de contato se possível, e continuar enxaguando
(P305+P351+P338); depois, procurar orientação médica.
\textbf{6.} Para os resíduos aquosos, após neutralização, e não para a
pia.
\textbf{7.} \ce{C7H6O3 + C4H6O3 -> C9H8O4 + C2H4O2}: C $7 + 4 = 9 + 2$,
H $6 + 6 = 8 + 4$, O $3 + 3 = 4 + 2$.
\textbf{8.} $2.00/138.0 = \qty{1.449e-2}{mol}$; no máximo
$\num{1.449e-2} \times 180.0 = \qty{2.61}{g}$ de aspirina.
\textbf{9.} O sólido bruto contém ácido salicílico que não reagiu, ácido
etanoico e água: sua massa não é a da aspirina, e sua pureza é
desconhecida.
\textbf{10.} O que fica dissolvido após o resfriamento se perde: uma
\omterm{def:b1:precipitation:solubility}{solubilidade} pequena a frio e grande a quente significa que o produto se
dissolve a quente e volta quase inteiro a frio.
\textbf{11.} Cada mililitro a mais mantém dissolvida a sua parte de
produto a frio.
\textbf{12.} O crescimento lento dá \omterm{def:b1:crystals-metals:crystal}{cristais} maiores e mais regulares, que
aprisionam menos água-mãe e impurezas; o gelo diminui a \omterm{def:b1:precipitation:solubility}{solubilidade} e,
portanto, a perda.
\textbf{13.} $\qty{0.045}{L} \times \qty{3.3}{g/L} = \qty{0.15}{g}$.
\textbf{14.} $1.95/2.61 = \qty{75}{\percent}$ (\qty{74.7}{\percent}).
\textbf{15.} $807/6 = \qty{134.5}{\celsius}$.
\textbf{16.} Desvios $-0.5$, $0.5$, $-1.5$, $0.5$, $-0.5$, $1.5$; soma
dos quadrados 5.5; $s = \sqrt{5.5/5} = \qty{1.05}{\celsius}$.
\textbf{17.} $u_A = 1.05/\sqrt6 = \qty{0.43}{\celsius}$.
\textbf{18.} Semiamplitude \qty{0.5}{\celsius}: $0.5/\sqrt3 = \qty{0.29}{\celsius}$.
\textbf{19.} $1/\sqrt3 = \qty{0.58}{\celsius}$;
$u = \sqrt{0.428^2 + 0.289^2 + 0.577^2} = \sqrt{0.600} = \qty{0.77}{\celsius}$.
\textbf{20.} $U = 2 \times 0.775 = \qty{1.5}{\celsius}$:
$T = (134.5 \pm 1.5)$~\unit{\celsius}, $k = 2$.
\textbf{21.} Ele funde baixo e ao longo de uma faixa de \qty{8}{\celsius}:
é impuro (ácido salicílico, ácido etanoico, água).
\textbf{22.} Dado na unidade, o valor fica a $\pm\qty{0.5}{\celsius}$:
$u_{\mathrm{ref}} = 0.5/\sqrt3 = \qty{0.29}{\celsius}$.
\textbf{23.} $2.3/6.0 = 0.38$ e $3.3/6.0 = 0.55$. O produto bruto continha
ácido salicílico (mesmo $R_f$ que a referência) e aspirina; após a
\omterm{def:b1:lab-techniques-1:recrystallisation}{recristalização}, só resta a mancha da aspirina: o ácido salicílico foi
removido, pelo menos abaixo do que a placa detecta.
\textbf{24.} A aspirina se decompõe ao fundir; aquecida lentamente, ela
tem tempo de se decompor, e a mistura formada funde mais baixo. O valor
tabelado se refere ao aquecimento rápido, como no banco.
\textbf{25.} $z = |134.5 - 135|/\sqrt{0.600 + 0.083} = 0.5/0.827 =$
\textbf{0.60}: bem abaixo de 2, os pontos de fusão medido e tabelado são
compatíveis.
\end{solution}
